% GENERATED FILE. Edit canonical lessons, then run npm run generate:books.
% canonical-source-hash: fnv1a64:9a972b2ef06aa779
% canonical-lessons: GE-C428-abfluss, GE-C428-verstopfen, GE-C428-tropfen, GE-C428-undicht, GE-C428-klempner

\chapter{Ärger im Bad}
\label{ch:ge-arger-im-bad}
\hlchaptermodality{428}

\begin{chapteropening}
\textbf{By the end of this chapter:} \emph{I can report a bathroom problem: the drain, a blockage, dripping, a leak and the plumber.}

Complete the last lesson of chapter 428: I can report a bathroom problem: the drain, a blockage, dripping, a leak and the plumber.
\end{chapteropening}

\section[der Abfluss]{der Abfluss — the drain, the plughole}

\label{lesson:GE-C428-abfluss}



\begin{quote}
Before the new one: say the German for to turn up, then the German for to shut off.
\end{quote}

\subsection*{You'll want to know: der Abfluss}
\textbf{der Abfluss} — \textquotedblleft{}the drain, the plughole\textquotedblright{}.

\begin{quote}
Im Abfluss der Dusche sind viele Haare.
\end{quote}

\begin{quote}
\emph{There is a lot of hair in the shower drain.}
\end{quote}

Plural: \emph{die Abflüsse}.

\begin{cousinweb}
\emph{Abfluss} is built from \emph{ab} 'away' and \emph{Fluss} 'flow, river', inherited from Old High German \emph{fluz} (EtymDB), from the verb for flowing.
\end{cousinweb}

\begin{grammarlens}[title={What you've built}]
One more word for telling, explaining and describing at B1.
\end{grammarlens}

\subsection*{Your turn}
\begin{itemize}
\raggedright
  \item \emph{Say it:} \emph{der Abfluss}
  \item \emph{Say it:} \emph{der Abfluss}, once more
  \item \emph{Say it:} a sentence of your own with \emph{der Abfluss}
  \item \emph{Recall:} say the German for to turn up, then the German for to shut off, then say \emph{der Abfluss} again
\end{itemize}

\begin{tcolorbox}[breakable,colback=teal!4,colframe=teal!35!black,title={Before you move on}]
What does der Abfluss mean? (\textquotedblleft{}The drain, the plughole\textquotedblright{}.) Say it once more.
\end{tcolorbox}
\section[verstopfen]{verstopfen — to block, to clog}

\label{lesson:GE-C428-verstopfen}



\begin{quote}
Before the new one: say the German for to shut off, then the German for the drain.
\end{quote}

\subsection*{You'll want to know: verstopfen}
\textbf{verstopfen} — \textquotedblleft{}to block, to clog\textquotedblright{}.

\begin{quote}
Die Haare verstopfen den Abfluss.
\end{quote}

\begin{quote}
\emph{The hair blocks the drain.}
\end{quote}

\emph{Die Toilette ist verstopft} — the toilet is blocked. \emph{Verstopft} also means constipated.

\begin{cousinweb}
\emph{verstopfen} is inherited from Middle High German \emph{verstophen} and Old High German \emph{firstopfōn} (EtymDB), built on \emph{stopfen} 'to plug, to stuff'.
\end{cousinweb}

\begin{grammarlens}[title={What you've built}]
One more word for telling, explaining and describing at B1.
\end{grammarlens}

\subsection*{Your turn}
\begin{itemize}
\raggedright
  \item \emph{Say it:} \emph{verstopfen}
  \item \emph{Say it:} \emph{verstopfen}, once more
  \item \emph{Say it:} a sentence of your own with \emph{verstopfen}
  \item \emph{Recall:} say the German for to shut off, then the German for the drain, then say \emph{verstopfen} again
\end{itemize}

\begin{tcolorbox}[breakable,colback=teal!4,colframe=teal!35!black,title={Before you move on}]
What does verstopfen mean? (\textquotedblleft{}To block, to clog\textquotedblright{}.) Say it once more.
\end{tcolorbox}
\section[tropfen]{tropfen — to drip}

\label{lesson:GE-C428-tropfen}



\begin{quote}
Before the new one: say the German for the drain, then the German for to block.
\end{quote}

\subsection*{You'll want to know: tropfen}
\textbf{tropfen} — \textquotedblleft{}to drip\textquotedblright{}.

\begin{quote}
Die ganze Nacht tropft das Wasser in der Dusche.
\end{quote}

\begin{quote}
\emph{The water drips in the shower all night.}
\end{quote}

\emph{Der Tropfen} — the drop. \emph{Es tropft} — it is dripping.

\begin{cousinweb}
\emph{tropfen} is inherited through Middle High German \emph{tropfen} from Old High German \emph{tropfōn} (EtymDB). English \emph{drop} is a cousin.
\end{cousinweb}

\begin{grammarlens}[title={What you've built}]
One more word for telling, explaining and describing at B1.
\end{grammarlens}

\subsection*{Your turn}
\begin{itemize}
\raggedright
  \item \emph{Say it:} \emph{tropfen}
  \item \emph{Say it:} \emph{tropfen}, once more
  \item \emph{Say it:} a sentence of your own with \emph{tropfen}
  \item \emph{Recall:} say the German for the drain, then the German for to block, then say \emph{tropfen} again
\end{itemize}

\begin{tcolorbox}[breakable,colback=teal!4,colframe=teal!35!black,title={Before you move on}]
What does tropfen mean? (\textquotedblleft{}To drip\textquotedblright{}.) Say it once more.
\end{tcolorbox}
\section[undicht]{undicht — leaky, not watertight}

\label{lesson:GE-C428-undicht}



\begin{quote}
Before the new one: say the German for to block, then the German for to drip.
\end{quote}

\subsection*{You'll want to know: undicht}
\textbf{undicht} — \textquotedblleft{}leaky, not watertight\textquotedblright{}.

\begin{quote}
Das Dach über dem Bad ist undicht.
\end{quote}

\begin{quote}
\emph{The roof over the bathroom leaks.}
\end{quote}

\emph{Dicht} means tight, sealed; \emph{undicht} is its opposite.

\begin{cousinweb}
\emph{undicht} is built from \emph{un-} 'not' and \emph{dicht} 'tight, dense', an inherited Germanic word (EtymDB lists \emph{dicht}).
\end{cousinweb}

\begin{grammarlens}[title={What you've built}]
One more word for telling, explaining and describing at B1.
\end{grammarlens}

\subsection*{Your turn}
\begin{itemize}
\raggedright
  \item \emph{Say it:} \emph{undicht}
  \item \emph{Say it:} \emph{undicht}, once more
  \item \emph{Say it:} a sentence of your own with \emph{undicht}
  \item \emph{Recall:} say the German for to block, then the German for to drip, then say \emph{undicht} again
\end{itemize}

\begin{tcolorbox}[breakable,colback=teal!4,colframe=teal!35!black,title={Before you move on}]
What does undicht mean? (\textquotedblleft{}Leaky, not watertight\textquotedblright{}.) Say it once more.
\end{tcolorbox}
\section[der Klempner]{der Klempner — the plumber}

\label{lesson:GE-C428-klempner}



\begin{quote}
Before the new one: say the German for to drip, then the German for leaky.
\end{quote}

\subsection*{You'll want to know: der Klempner}
\textbf{der Klempner} — \textquotedblleft{}the plumber\textquotedblright{}.

\begin{quote}
Der Klempner kommt erst am Montag.
\end{quote}

\begin{quote}
\emph{The plumber is not coming until Monday.}
\end{quote}

Same form in the plural. In Austria: \emph{der Installateur}.

\begin{cousinweb}
\emph{Klempner} first meant a tinsmith, someone who hammers sheet metal. It is usually traced to \emph{klempern} 'to bang on metal'; EtymDB lists the word without a source, so the origin is hedged.
\end{cousinweb}

\begin{grammarlens}[title={What you've built}]
One more word for telling, explaining and describing at B1.
\end{grammarlens}

\subsection*{Your turn}
\begin{itemize}
\raggedright
  \item \emph{Say it:} \emph{der Klempner}
  \item \emph{Say it:} \emph{der Klempner}, once more
  \item \emph{Say it:} a sentence of your own with \emph{der Klempner}
  \item \emph{Recall:} say the German for to drip, then the German for leaky, then say \emph{der Klempner} again
\end{itemize}

\begin{tcolorbox}[breakable,colback=teal!4,colframe=teal!35!black,title={Before you move on}]
What does der Klempner mean? (\textquotedblleft{}The plumber\textquotedblright{}.) Say it once more.
\end{tcolorbox}
